\documentclass[11pt]{article}

\usepackage[margin=1in]{geometry}
\usepackage{amsmath,amssymb}
\usepackage{booktabs}
\usepackage[colorlinks=true,linkcolor=blue,urlcolor=blue]{hyperref}

\newcommand{\eps}{\varepsilon}
\newcommand{\dd}{\mathrm{d}}
\newcommand{\Li}{\mathrm{Li}}

\title{Solving the NLO phase-space master integrals\\
by differential equations\\[1ex]
\large Method note for the $pp \to h X$ workflow}
\author{Max Zhang}
\date{\today}

\begin{document}
\maketitle

\begin{abstract}
This note explains, on the concrete example of the seven NLO
double-parton phase-space masters of the unpolarized $pp \to hX$
channel, how the differential-equation (DE) method produces exact
master integrals: how the system of first-order equations is derived
from integration-by-parts identities, how boundary conditions enter,
and --- the point most relevant for single-inclusive kinematics ---
how the method reconstructs the \emph{exact} endpoint exponents
$(1-v-w)^{\,n+a\eps}$ that the plus-distribution expansion of the
partonic cross section requires, even though every explicit result is
only ever written as an expansion in $\eps$.
\end{abstract}

\section{The objects}
\label{sec:objects}

At NLO the real correction to $ab \to c\,X$ is a $2\to3$ process with
the fragmenting parton momentum $k_c$ fixed and two unobserved
partons integrated over their on-shell phase space.  By reverse
unitarity the phase-space integrals are loop integrals with two cut
propagators: after eliminating one unobserved momentum by momentum
conservation, every integral is of the form
\begin{equation}
  I_{\nu_1\nu_2\nu_3\nu_4}
  \;=\;
  \int \dd^d k_e\;
  \frac{\delta^{+}\!\big(k_e^2\big)\,
        \delta^{+}\!\big((P-k_e)^2\big)}
       {D_2^{\nu_2}\, D_3^{\nu_3}},
  \qquad P = k_a + k_b - k_c ,
  \label{eq:integral}
\end{equation}
with massless external momenta $k_a, k_b, k_c$ and family-dependent
propagators $D_{2,3}$ (e.g.\ $(k_e+k_c)^2$, $(k_e-k_a)^2$, \dots).
The integrals depend on the three invariants only through an overall
power of $s_{12} = 2k_a\!\cdot\!k_b$ and the two dimensionless ratios
\begin{equation}
  v = \frac{2k_a\!\cdot\!k_c}{2k_a\!\cdot\!k_b},
  \qquad
  w = \frac{2k_b\!\cdot\!k_c}{2k_a\!\cdot\!k_b},
  \qquad
  0 < v,\; 0 < w,\; v + w < 1 ,
\end{equation}
where $1 - v - w = P^2/s_{12}$ is the normalized invariant mass of
the unobserved pair: the boundary $v+w \to 1$ is the soft endpoint
where the plus-distributions of the collinear-factorized cross
section live.  The IBP reduction of the full amplitude (Kira, with
the cut declared so that any sector with a vanished cut propagator is
zero) leaves seven master integrals: the pure two-body phase-space
volume $I_{\rm vol}$ and six one-and-two-propagator integrals, one
per family.

\section{Deriving the differential equations}
\label{sec:derive}

The masters are functions of the invariants, so their derivatives are
again integrals of the same families.  Concretely, apply the three
Euler operators $k\!\cdot\!\partial_k$ ($k = k_a, k_b, k_c$) under
the integral sign in Eq.~\eqref{eq:integral}: each operator acts on
the propagators (raising powers and producing scalar products), and
the loop-momentum scalar products are rewritten back through the
propagators of the family --- the families are complete, so this
rewriting is a linear solve.  On the other side, acting on
$I(s_{12}, s_{13}, s_{23})$ as a function, the same operators give
\begin{equation}
  k_a\!\cdot\!\partial_{k_a} = s_{12}\partial_{12} + s_{13}\partial_{13},
  \quad
  k_b\!\cdot\!\partial_{k_b} = s_{12}\partial_{12} + s_{23}\partial_{23},
  \quad
  k_c\!\cdot\!\partial_{k_c} = s_{13}\partial_{13} + s_{23}\partial_{23},
\end{equation}
an invertible linear system for the invariant derivatives.  The
result is a set of integrals with indices raised or lowered by one
unit; a small dedicated IBP reduction (the derivative integrals raise
\emph{cut} indices, which the amplitude reduction never needed) maps
them back onto the masters.  The output is a closed linear system
\begin{equation}
  \partial_v \vec I \;=\; A_v(v,w;\eps)\, \vec I,
  \qquad
  \partial_w \vec I \;=\; A_w(v,w;\eps)\, \vec I,
  \label{eq:system}
\end{equation}
with $A_{v,w}$ \emph{exact rational functions} of $v,w,\eps$.  Two
structural checks certify the construction:
\begin{itemize}
\item \textbf{Flatness.}  The integrability condition
  $\partial_w A_v - \partial_v A_w + [A_v, A_w] = 0$ must hold
  identically; for our system it does, exactly.  This single check
  exercises the Euler construction, the completeness rewriting, and
  the IBP tables at once.
\item \textbf{Alphabet.}  The matrix entries only have poles on the
  hypersurfaces
  \begin{equation}
    \{\, v,\; w,\; 1-v,\; 1-w,\; v+w,\; 1-v-w \,\},
    \label{eq:alphabet}
  \end{equation}
  the \emph{symbol alphabet} of the problem: every singularity of
  every master lies on one of these six lines.
\end{itemize}
For the NLO system the matrices come out lower-triangular in a
convenient ordering: the volume decouples,
$\partial_v \log I_{\rm vol} = -\eps\,\partial_v\log(1-v-w)$
(so $I_{\rm vol} \propto (1-v-w)^{-\eps}$, as expected from
$(P^2)^{d/2-2}$), and each remaining master couples only to itself
and to $I_{\rm vol}$: six scalar first-order equations with a known
source.

\section{Solving the system}
\label{sec:solve}

\subsection{Exact integrating factors}

The diagonal entries of $A_{v,w}$ are $\dd\log$-forms in the alphabet
with residues linear in $\eps$, so each master has an exact
integrating factor: a product $\prod_l \alpha_l^{\rho_l(\eps)}$ over
the letters.  Writing $I_i = p_i(v,w;\eps)\, h_i(v,w;\eps)$ with the
prefactor $p_i$ chosen to absorb the diagonal, the remaining
equations are pure quadratures,
\begin{equation}
  \partial_v h_i = p_i^{-1} (A_v)_{i1}\, I_{\rm vol},
  \qquad
  \partial_w h_i = p_i^{-1} (A_w)_{i1}\, I_{\rm vol},
\end{equation}
whose right-hand sides are rational functions times products of
letters with $\eps$-linear exponents.  Expanding those exponents
order by order in $\eps$ produces integrands of the form
(rational)$\times$(logarithms of letters), and each order integrates
in closed form: the solutions live in the space of classical
polylogarithms of the letters, $\log$, $\Li_2$, $\Li_3$, \dots\ ---
at weight $\le 3$ for our purposes.  (In the general, non-triangular
case one first rotates the basis to Henn's canonical form
$A = \eps \sum_l M_l\, \dd\log \alpha_l$ and solves by iterated
integrals; the triangular NLO system lets us shortcut that machinery
while keeping its essential feature, the exact $\dd\log$ prefactors.)

\subsection{Boundary conditions}

A first-order system in two variables needs one constant of
integration per master (as a function of $\eps$).  These come from
integrals that are simpler than the masters themselves:
\begin{itemize}
\item the two-body phase-space volume is elementary,
  \begin{equation}
    I_{\rm vol}
    = \frac{2\pi^{\frac{d-1}{2}}}{\Gamma\!\big(\tfrac{d-1}{2}\big)}\,
      2^{1-d}\, \big(s_{12}(1{-}v{-}w)\big)^{\frac{d-4}{2}} ,
    \qquad d = 4-2\eps ;
  \end{equation}
\item the remaining constants are fixed by regularity conditions
  (a master must not be singular on alphabet lines where its
  definition is manifestly finite) and, where needed, by matching a
  single independent evaluation --- here a two-angle numerical
  integration of Eq.~\eqref{eq:integral} in the $P$ rest frame,
  performed at negative $\eps$ where the angular integrals converge
  absolutely, and at high precision so that the $\eps$-series
  coefficients can be identified exactly.
\end{itemize}
The crucial point, developed next, is \emph{how little} boundary
information is needed: only $\eps$-\emph{expanded} numbers, never
exact-in-$\eps$ functions.

\section{Endpoint singularities from the DE:\\
how $(1-v-w)^{-1-a\eps}$ is reconstructed}
\label{sec:endpoint}

The partonic cross section must ultimately be expanded in
distributions at the soft endpoint,
\begin{equation}
  (1-v-w)^{-1-a\eps}
  \;=\;
  -\frac{\delta(1-v-w)}{a\eps}
  + \left[\frac{1}{1-v-w}\right]_+
  - a\eps \left[\frac{\ln(1-v-w)}{1-v-w}\right]_+
  + \dots,
  \label{eq:plus}
\end{equation}
and this expansion needs the exponent $-1-a\eps$ \emph{exactly}: an
$\eps$-truncated $\frac{1}{1-v-w}\big(1 - a\eps\ln(1-v-w)+\dots\big)$
carries no $\delta$-term.  At first sight a solution obtained order
by order in $\eps$ seems to have discarded exactly this information.
It has not, for the following reason.

The connection matrices in Eq.~\eqref{eq:system} were never expanded:
they are exact rational functions of $\eps$.  Near the soft surface
the system has a \emph{regular singular point} in the sense of
Fuchsian ODE theory,
\begin{equation}
  A_w \;=\; \frac{R(\eps)}{w - (1-v)} + \mathcal{O}(1),
\end{equation}
and the local solution space is governed entirely by the residue
matrix $R(\eps)$: by the Frobenius construction,
\begin{equation}
  \vec I(v,w)
  \;=\;
  \sum_i c_i(\eps)\,(1-v-w)^{\lambda_i(\eps)}\,
  \vec h_i(v,w;\eps),
  \qquad
  \lambda_i(\eps) \in \operatorname{eig} R(\eps),
  \label{eq:frobenius}
\end{equation}
with $\vec h_i$ analytic on the surface.  The exponents
$\lambda_i(\eps) = n_i + a_i\eps$ are \emph{eigenvalues of an exact
rational matrix} --- they are exact in $\eps$ by construction,
whatever expansion is used later for the analytic parts.  In our
system one reads off, for example,
\begin{equation}
  I_{\rm vol} \sim (1-v-w)^{-\eps},
  \qquad
  I_{3} \sim (1-v-w)^{-1-2\eps} \times \big(\dots\big)
  \;+\; (1-v-w)^{-\eps} \times \big(\dots\big),
\end{equation}
the $-1-2\eps$ being precisely the residue eigenvalue $d-5$ of the
corresponding diagonal entry: this is the seed of
Eq.~\eqref{eq:plus}.

The practical consequences:
\begin{enumerate}
\item \textbf{Solve per branch.}  The solution is organized as
  (exact endpoint prefactor) $\times$ ($\eps$-expanded analytic
  part), Eq.~\eqref{eq:frobenius}.  Because the non-analyticity is
  factored out exactly, the $\eps$-expansion of the analytic part
  commutes with the endpoint limit: nothing is lost by truncating.
\item \textbf{Boundary data can be $\eps$-expanded.}  The constants
  $c_i(\eps)$ multiply the branches; they do not shape them.  This is
  the resolution of the apparent tension above: the singular
  structure comes from the exact DE, the boundary only normalizes
  each branch --- and a normalization is needed only as a series.
  (One caveat: the branch feeding the $\delta$-term in
  Eq.~\eqref{eq:plus} carries an explicit $1/\eps$, so its constant
  is needed to one order in $\eps$ beyond the target accuracy.)
\item \textbf{Store masters branch-factored.}  A naive expanded
  series with $\ln^k(1-v-w)$ terms is valid pointwise in the open
  region but meaningless as a distribution across the endpoint.  The
  artifact format for solved masters is therefore
  $(1-v-w)^{n+a\eps} \times$ (analytic $\eps$-series): this is the
  form the plus-distribution extraction against the declared
  prefactor $f_1(x_a) f_1(x_b) D_1(z_h)$ consumes directly.  The same
  discipline applies at NNLO, where the double-real masters carry
  branches like $(1-v-w)^{-1-4\eps}$ from two unresolved partons.
\end{enumerate}

\section{Validation}
\label{sec:validation}

Independent checks, in increasing strength:
\begin{itemize}
\item the exact flatness of the connection
  (Section~\ref{sec:derive});
\item the volume row reproduces the known closed form exactly;
\item every master is evaluated numerically at interior phase-space
  points by the direct two-angle representation of
  Eq.~\eqref{eq:integral} at several negative values of $\eps$, and
  compared with the DE solution: agreement at one point fixes the
  constant, agreement at a second point (and at multiple $\eps$) is a
  genuine test of the full function.
\end{itemize}

\section{Outlook: the same method at NNLO}
\label{sec:outlook}

Nothing in Sections~\ref{sec:derive}--\ref{sec:endpoint} is specific
to NLO.  For the $347$ double-real NNLO masters the same Euler
construction gives the system, the same flatness identity certifies
it, and the same Frobenius analysis at the soft surface delivers the
exact endpoint exponents for the NNLO plus-distributions.  The
differences are quantitative: the alphabet grows, the systems no
longer arrange themselves triangularly (a genuine rotation to
canonical $\eps$-form is then the right tool), boundary conditions
are fixed corner by corner in the $(v,w)$ square, and the function
space reaches weight~4.  The NLO system solved here is the
methodological template --- and its solution enters the NNLO
subtraction structure directly, since the NLO masters convolved with
splitting kernels build the collinear counterterms of the NNLO hard
function.

\end{document}
